
\subsection{Conclusiones}

El análisis realizado permite caracterizar la evolución de ROS como una
transición desde una arquitectura de comunicaciones propia y parcialmente
centralizada hacia una arquitectura distribuida sustentada en estándares
industriales. En ROS~1, el proceso maestro concentra el registro y la
resolución de nombres, lo que lo convierte en un punto único de falla y
condiciona el despliegue en redes extensas o poco confiables. ROS~2 elimina
esta dependencia mediante el descubrimiento distribuido que provee RTPS.

La adopción de DDS y RTPS inscribe a ROS~2 en un marco normativo consolidado,
definido por la OMG en las capas de \textit{middleware} e interoperabilidad y
por la IETF en las capas de transporte y red. Esta organización en capas
asigna a cada nivel una función definida: DDS establece el modelo de datos,
RTPS el protocolo de interoperabilidad, UDP el transporte e IP el
encaminamiento. La elección de UDP como transporte resulta coherente con este
diseño, dado que habilita el descubrimiento mediante \textit{multicast} y
traslada a RTPS la implementación de la confiabilidad, que se aplica solo
cuando la política de calidad de servicio lo requiere.

Desde la perspectiva de los sistemas de comunicaciones, las políticas de QoS
constituyen el aporte más significativo de ROS~2, ya que permiten adecuar el
comportamiento de cada flujo de datos a las características del enlace y de la
aplicación, por ejemplo, priorizando la actualidad de los datos de sensores
sobre la garantía de entrega. Como contrapartida, incorporan un factor de
complejidad en el diseño: los perfiles incompatibles impiden establecer la
comunicación, por lo que su configuración debe definirse de manera coherente
en todo el sistema.

La capa de abstracción RMW desacopla las aplicaciones de la implementación del
\textit{middleware} y habilita la coexistencia de múltiples proveedores con
modelos de licenciamiento diferenciados. La incorporación de Zenoh como
implementación de primer nivel demuestra que esta abstracción admite incluso
alternativas no basadas en DDS, orientadas a escenarios en los que el
descubrimiento mediante \textit{multicast} presenta limitaciones, como las
redes inalámbricas.

Por último, ROS~2 amplía su alcance hacia sistemas con requisitos de tiempo
real, mediante el soporte de núcleos Linux con PREEMPT\_RT, y hacia
microcontroladores, mediante micro-ROS y el estándar DDS-XRCE. No obstante, el
determinismo temporal depende del conjunto de la pila de ejecución y no
únicamente del \textit{middleware}. Asimismo, la vinculación de cada
distribución con una versión LTS de Ubuntu constituye un factor a considerar
en la planificación del ciclo de vida de los sistemas desplegados.